\chapter*{CONCLUSION ET PERSPECTIVES}
\markboth{\MakeUppercase{CONCLUSION GÉNÉRALE}}{}
\addcontentsline{toc}{chapter}{CONCLUSION GÉNÉRALE}
\adjustmtc
\thispagestyle{MyStyle}

% Guide pour la rédaction de la Conclusion et Perspectives

La \textbf{Conclusion Générale} a pour objectif de résumer les contributions du projet et d’ouvrir des perspectives d’évolution. Elle doit inclure les éléments suivants :

\begin{itemize}
    \item \textbf{Synthèse des réalisations :} Récapitulez les principales étapes du projet et les résultats obtenus.
    \item \textbf{Impact et apports :} Mettez en avant l'importance du projet et les bénéfices qu’il apporte aux utilisateurs ou à l’entreprise.
    \item \textbf{Limites :} Identifiez brièvement les contraintes ou améliorations possibles.
    \item \textbf{Perspectives d’évolution :} Décrivez les pistes d’amélioration et d’extension du projet :
    \begin{itemize}
        \item \textbf{À court terme :} Actions immédiates après la mise en place du projet.
        \item \textbf{À moyen terme :} Évolutions possibles dans les mois ou années à venir.
        \item \textbf{À long terme :} Vision future et extensions potentielles du projet.
    \end{itemize}
\end{itemize}

\noindent Cette conclusion doit être rédigée de manière synthétique et mettre en avant la valeur ajoutée du projet ainsi que son potentiel d’évolution.
